\subsection*{暂时不画函数曲线，改画两条数轴}
把输入数轴上的小点送到输出数轴，函数就成为点的移动规则。在 $a$ 附近放两个很近的点，它们的输出间距约是输入间距的 $f'(a)$ 倍。正导数保序，负导数反序，绝对值表示局部放大程度。

对平方函数，在三附近，输入间距大约被放大六倍；在负三附近，倍率是负六，输出次序反转。零点的导数为零，说明一阶间距被压掉，但二阶变化仍在，所以附近函数没有变成常数。

\subsection*{线性代数与微积分在这里接上}
局部表达 $f(a+h)=f(a)+f'(a)h+o(h)$，把固定输出位置与输入微扰的线性响应分开。推广到多变量，数字 $f'(a)$ 会变成一个线性映射或雅可比矩阵；它描述附近小箭头怎样改变，正如前面的矩阵变换。

复合函数先通过内层局部变换，再通过外层局部变换，所以导数相乘。多变量情形使用矩阵乘法，顺序依旧是右边先做。链式法则因此不仅是一条求导技巧，也是一条局部动作的复合规律。

请在纸上画输入与输出两条平行数轴，用连线表示三个邻近输入的像，再根据间距预测导数正负。这个视角尤其适合解释函数图像难以直接画出的高维问题，也为神经网络中的敏感性分析准备了语言。
